%!TEX TS-program = pdflatex
\documentclass[aspectratio=169,10pt]{beamer}
\usepackage[T2A]{fontenc}
\usepackage[utf8]{inputenc}
\usepackage[english,russian]{babel}
\usepackage{cmap}
\usepackage{amsmath,amssymb,mathtools,booktabs,array,etoolbox,PTSans,graphicx}
\graphicspath{{figures/}}
\renewcommand{\familydefault}{\sfdefault}
\renewcommand{\rmdefault}{\sfdefault}
\newbool{russian}\booltrue{russian}
\usepackage{HSE-theme/beamerthemeHSE}
\hypersetup{bookmarks=false}
\setbeamertemplate{blocks}[rounded][shadow=false]
\setbeamerfont{frametitle}{size=\Large}
\setbeamerfont{framesubtitle}{size=\small}
\setbeamerfont{block title}{size=\normalsize}
\setbeamerfont{block body}{size=\small}
\setlength{\abovedisplayskip}{6pt}
\setlength{\belowdisplayskip}{6pt}
\setlength{\abovedisplayshortskip}{4pt}
\setlength{\belowdisplayshortskip}{4pt}
\newcommand{\R}{\mathbb{R}}
\newcommand{\figtall}[1]{\begin{center}\includegraphics[height=.50\textheight,keepaspectratio]{#1}\end{center}}
\newcommand{\figshort}[1]{\begin{center}\includegraphics[height=.38\textheight,keepaspectratio]{#1}\end{center}}
\title[Методы оптимизации]{Методы оптимизации\\в машинном обучении}
\subtitle{Лекция 8. Выбор шага: точный поиск, условие Армихо и гарантии}
\author[Н. Лукьяненко]{Никита Лукьяненко}
\institute{Совместная программа СМОЛ ГУ и ФКН НИУ ВШЭ}
\date{2026/27 учебный год}
\begin{document}
% Slide 1
\frame[plain]{\titlepage}

% Slide 2
\begin{frame}{Шаг и направление — две разные величины}
\small
Общий градиентный метод задаёт итерацию
\[
\boxed{x_{k+1}=x_k+\alpha_k d_k,\qquad \alpha_k>0.}
\]
Здесь $d_k$ отвечает за направление, а $\alpha_k$ --- за длину перемещения.\\[.5em]
Для градиентного спуска $d_k=-\nabla f(x_k)$, поэтому
\[
x_{k+1}=x_k-\alpha_k\nabla f(x_k).
\]
\begin{block}{Постановка вопроса}
Как выбрать $\alpha_k$, если градиент уже вычислен?
\end{block}
\end{frame}

% Slide 3
\begin{frame}{Одна функция, три режима выбора шага}
\small
\figtall{step_1d.pdf}
\vspace{-.3em}
Для $f(x)=x^2$ при постоянном шаге:
\[
x_{k+1}=(1-2\alpha)x_k.
\]
Характер поведения целиком определяется числом $|1-2\alpha|$.
\end{frame}

% Slide 4
\begin{frame}{В многомерном случае важен спектр}
\small
\figtall{step_contours.pdf}
\vspace{-.4em}
Та же квадратичная задача, но разные $\alpha$: медленное движение, затухающие колебания, расходимость.
\end{frame}

% Slide 5
\begin{frame}{Известный результат для квадратичной функции}
\small
Пусть $f(x)=\frac12x^\top Ax-b^\top x$ и $A=A^\top\succ0$.
Тогда градиентный спуск с постоянным шагом сходится из любой начальной точки тогда и только тогда, когда
\[
\boxed{0<\alpha<\frac2L,\qquad L=\lambda_{\max}(A).}
\]
\figshort{spectral_threshold.pdf}
\end{frame}

% Slide 6
\begin{frame}{Почему фиксированного шага недостаточно}
\small
Для общей гладкой функции константа $L$ может быть неизвестна, а её глобальная оценка --- чрезмерно большой.
\begin{block}{Новая идея}
Вместо выбора шага заранее попробуем искать его \textbf{в самой текущей точке}, пользуясь значениями $f$ и $\nabla f$.
\end{block}
Нужно формально определить, что считать хорошим шагом, и доказать, что его можно найти.
\end{frame}

% Slide 7
\begin{frame}{Многомерная задача превращается в одномерную}
\small
Зафиксируем точку $x\in\R^n$ и ненулевое направление $d\in\R^n$.
\begin{block}{Определение}
Ограничение функции на прямую $x+\alpha d$:
\[
\boxed{\varphi(\alpha)=f(x+\alpha d),\quad \alpha\ge0.}
\]
\end{block}
Теперь выбор шага --- одномерная задача по переменной $\alpha$.
\end{frame}

% Slide 8
\begin{frame}{Первая производная вдоль направления}
\small
\begin{block}{Утверждение 1}
Если $f$ дифференцируема, то
\[
\boxed{\varphi'(\alpha)=\nabla f(x+\alpha d)^\top d.}
\]
\end{block}
\textbf{Доказательство.} Положим $z(\alpha)=x+\alpha d$. Тогда по правилу дифференцирования сложной функции
\[
\varphi'(\alpha)=\nabla f(z(\alpha))^\top z'(\alpha)
=\nabla f(x+\alpha d)^\top d.\quad\square
\]
\end{frame}

% Slide 9
\begin{frame}{Вторая производная вдоль направления}
\small
\begin{block}{Утверждение 2}
Если $f\in C^2$, то
\[
\boxed{\varphi''(\alpha)=d^\top\nabla^2 f(x+\alpha d)d.}
\]
\end{block}
\textbf{Доказательство.} Дифференцируем $\nabla f(x+\alpha d)^\top d$ по $\alpha$. Вектор $d$ постоянен:
\[
\frac{d}{d\alpha}\nabla f(x+\alpha d)=\nabla^2f(x+\alpha d)d.
\]
Подставляя, получаем формулу. $\square$
\end{frame}

% Slide 10
\begin{frame}{Частный случай: квадратичная функция}
\small
Пусть $f(z)=\frac12z^\top Az-b^\top z+c$, $A=A^\top$.
Для $\varphi(\alpha)=f(x+\alpha d)$ получаем \textbf{точное} разложение:
\[
\varphi(\alpha)=f(x)+\alpha\nabla f(x)^\top d
+\frac{\alpha^2}{2}d^\top Ad.
\]
Следовательно,
\[
\varphi'(\alpha)=\nabla f(x)^\top d+\alpha d^\top Ad,
\qquad
\varphi''(\alpha)=d^\top Ad.
\]
Это соединяет квадратичные формы с одномерным поиском шага.
\end{frame}

% Slide 11
\begin{frame}{Направление убывания: общий случай}
\small
\begin{block}{Определение}
Ненулевой вектор $d$ называется \textbf{направлением убывания} функции $f$ в точке $x$, если
\[
\boxed{\nabla f(x)^\top d<0.}
\]
\end{block}
Тогда $\varphi'(0)<0$, и малые положительные шаги уменьшают $f$.\newline
\figshort{directional_slope.pdf}
\end{frame}

% Slide 12
\begin{frame}{Почему малый положительный шаг уменьшает функцию}
\small
По определению производной
\[
\varphi'(0)=\lim_{\alpha\to0+}\frac{\varphi(\alpha)-\varphi(0)}{\alpha}<0.
\]
Поэтому для достаточно малых $\alpha>0$
\[
\frac{\varphi(\alpha)-\varphi(0)}{\alpha}<0.
\]
Так как $\alpha>0$, отсюда
\[
\boxed{f(x+\alpha d)<f(x).}\quad\square
\]
\textit{Важно:} здесь мы пока доказали убывание, но не нашли величину шага.
\end{frame}

% Slide 13
\begin{frame}{Антиградиент действительно убывает}
\small
Для $d=-\nabla f(x)$ при $\nabla f(x)\neq0$
\[
\nabla f(x)^\top d=-\|\nabla f(x)\|_2^2<0.
\]
Более того, среди единичных направлений $\|d\|_2=1$
\[
\min_{\|d\|_2=1}\nabla f(x)^\top d=-\|\nabla f(x)\|_2.
\]
\textbf{Почему?} По неравенству Коши--Буняковского
$\nabla f(x)^\top d\ge-\|\nabla f(x)\|_2$;
равенство достигается при $d=-\nabla f(x)/\|\nabla f(x)\|_2$.
\end{frame}

% Slide 14
\begin{frame}{Геометрия ограничения на прямую}
\small
\figtall{line_section_3d.pdf}
\vspace{-.4em}
Слева значения $f$ вдоль прямой $x+\alpha d$, справа --- та же зависимость как функция одной переменной $\varphi(\alpha)$.
\end{frame}

% Slide 15
\begin{frame}{Условие стационарности по шагу}
\small
Если $\varphi$ достигает \textbf{локального минимума} во внутренней точке $\alpha^*>0$, то по теореме Ферма
\[
\boxed{\varphi'(\alpha^*)=0.}
\]
Отсюда
\[
\nabla f(x+\alpha^*d)^\top d=0.
\]
\begin{alertblock}{Оговорка}
При граничном минимуме $\alpha^*=0$ равенство производной нулю не обязательно. И в общем случае стационарность ещё не доказывает минимальность.
\end{alertblock}
\end{frame}

% Slide 16
\begin{frame}{Когда стационарная точка шага — минимум}
\small
Если $\varphi''(\alpha)>0$ на рассматриваемом промежутке, то $\varphi$ строго выпукла.\\[.4em]
Поэтому у неё \textbf{не более одного} минимума, а любая внутренняя стационарная точка --- глобальный минимум $\varphi$.
\begin{block}{Для квадратичной функции}
Если $A\succ0$ и $d\ne0$, то
\[
\varphi''(\alpha)=d^\top Ad>0.
\]
Значит, одномерная задача имеет строгую выпуклую геометрию.
\end{block}
\end{frame}

% Slide 17
\begin{frame}{Точный поиск шага: определение}
\small
\begin{block}{Определение}
При фиксированных $x$ и $d$ \textbf{точным поиском шага} называют решение
\[
\boxed{\alpha^*\in\arg\min_{\alpha\ge0} f(x+\alpha d).}
\]
\end{block}
Это не оптимизация по всем координатам: меняется только $\alpha$.\\[.5em]
Если минимум достигается внутри полуоси, используем
$\varphi'(\alpha^*)=0$.
\end{frame}

% Slide 18
\begin{frame}{Внутренний минимум и граница допустимых шагов}
\small
Для дифференцируемой $\varphi$ минимум по $\alpha\ge0$ может быть
\begin{itemize}
\item внутренним ($\alpha^*>0$), тогда $\varphi'(\alpha^*)=0$;
\item граничным ($\alpha^*=0$), тогда достаточно $\varphi'(0)\ge0$ при выпуклой $\varphi$.
\end{itemize}
В дальнейшем берём $d=-g$, $g=\nabla f(x)\ne0$: тогда
\[
\varphi'(0)=-\|g\|_2^2<0,
\]
и нулевой шаг не может быть оптимальным.
\end{frame}

% Slide 19
\begin{frame}{Квадратичная функция: разложение вдоль антиградиента}
\small
Для $f(x)=\frac12x^\top Ax-b^\top x+c$, $A\succ0$, обозначим $g=Ax-b\ne0$.\\[.5em]
Вдоль $d=-g$:
\[
\boxed{\varphi(\alpha)=f(x)-\alpha\|g\|_2^2
+\frac{\alpha^2}{2}g^\top Ag.}
\]
Поскольку $g^\top Ag>0$, получена \textbf{строго выпуклая квадратичная} функция шага.
\end{frame}

% Slide 20
\begin{frame}{Точный шаг на квадратичной функции}
\small
Дифференцируем предыдущую формулу:
\[
\varphi'(\alpha)=-\|g\|_2^2+\alpha g^\top Ag.
\]
Из $\varphi'(\alpha^*)=0$ имеем
\[
\boxed{\alpha^*=\frac{\|g\|_2^2}{g^\top Ag}.}
\]
А $\varphi''(\alpha)=g^\top Ag>0$ доказывает, что это \textbf{единственный глобальный минимум} вдоль направления. $\square$
\end{frame}

% Slide 21
\begin{frame}{Насколько велик точный шаг?}
\small
\begin{block}{Утверждение 3: спектральные границы}
Если $A=A^\top\succ0$, $g\ne0$, $\mu=\lambda_{\min}(A)$, $L=\lambda_{\max}(A)$, то
\[
\boxed{\frac1L\le\alpha^*\le\frac1\mu.}
\]
\end{block}
\textbf{Доказательство.} По спектральным границам
\[
\mu\|g\|_2^2\le g^\top Ag\le L\|g\|_2^2.
\]
Делим $\|g\|_2^2>0$ на части неравенства, получая нужные оценки. $\square$
\end{frame}

% Slide 22
\begin{frame}{Пример спектральных границ}
\small
Пусть $A=\mathrm{diag}(1,9)$ и $g=(2,1)^\top$. Тогда
\[
\|g\|_2^2=5,\qquad g^\top Ag=4+9=13.
\]
Следовательно,
\[
\boxed{\alpha^*=\frac5{13}\approx0.385,\qquad
\frac19\le\frac5{13}\le1.}
\]
\figshort{rayleigh_steps.pdf}
\end{frame}

% Slide 23
\begin{frame}{Почему последовательные градиенты ортогональны}
\small
\begin{block}{Утверждение 4 (общий дифференцируемый случай)}
Пусть $g_k=\nabla f(x_k)\neq0$ и $\alpha_k^*>0$ --- точный минимум
$f(x_k-\alpha g_k)$ по $\alpha\ge0$. Тогда
\[
\boxed{g_{k+1}^{\top}g_k=0,\qquad g_{k+1}=\nabla f(x_{k+1}).}
\]
\end{block}
\textbf{Доказательство.} Зафиксируем $x_k$ и $g_k$ и положим
\[
\varphi_k(\alpha)=f(x_k-\alpha g_k).
\]
По правилу дифференцирования сложной функции
\[
\varphi_k'(\alpha)=-\nabla f(x_k-\alpha g_k)^\top g_k.
\]
Точный положительный шаг $\alpha_k^*$ --- внутренний минимум, поэтому
$\varphi_k'(\alpha_k^*)=0$. При этом $x_{k+1}=x_k-\alpha_k^*g_k$.
Значит,
\[
0=\varphi_k'(\alpha_k^*)
=-\nabla f(x_{k+1})^\top g_k=-g_{k+1}^\top g_k.\quad\square
\]
{\footnotesize\textit{Смысл: после точной минимизации вдоль $-g_k$ производная
по этому направлению равна нулю. Квадратичность $f$ не требуется.}}
\end{frame}

% Slide 24
\begin{frame}{Почему точный поиск даёт «зигзаг»}
\small
\figtall{orthogonal_gradients.pdf}
\vspace{-.4em}
На рисунке угол между $g_k$ и $g_{k+1}$ прямой, но это \textbf{не означает}, что каждый шаг ведёт напрямую к минимуму.
\end{frame}

% Slide 25
\begin{frame}{Сравнение нескольких точных шагов}
\small
\figtall{exact_path.pdf}
\vspace{-.4em}
Вытянутые линии уровня делают многократные повороты направления естественными даже при оптимальном шаге вдоль каждого антиградиента.
\end{frame}

% Slide 26
\begin{frame}{Цена точного поиска и следующая идея}
\small
Для квадратичной функции мы получили формулу $\alpha^*$. Но для общей $f$ задача
\[
\min_{\alpha\ge0} f(x+\alpha d)
\]
сама может потребовать множества вычислений функции или её производных.\\[.5em]
\begin{block}{Компромисс}
Не нужно находить \textbf{идеальный} шаг. Достаточно найти шаг с \textbf{доказуемым уменьшением функции}.
\end{block}
Это приводит к условию Армихо.
\end{frame}

% Slide 27
\begin{frame}{Линейное приближение вдоль направления}
\small
При малом $\alpha$ по дифференцируемости
\[
f(x+\alpha d)=f(x)+\alpha\nabla f(x)^\top d+o(\alpha).
\]
Обозначим \textbf{линейное приближение}
\[
\boxed{m_x(\alpha)=f(x)+\alpha\nabla f(x)^\top d.}
\]
Для направления убывания $\nabla f(x)^\top d<0$ оно задаёт отрицательный наклон.
\end{frame}

% Slide 28
\begin{frame}{Две величины уменьшения функции}
\small
\textbf{Фактическое уменьшение:}
\[
\Delta_{\text{факт}}(\alpha)=f(x)-f(x+\alpha d).
\]
\textbf{Уменьшение по линейному приближению:}
\[
\Delta_{\text{лин}}(\alpha)=f(x)-m_x(\alpha)
=-\alpha\nabla f(x)^\top d>0.
\]
Их не нужно путать: первая вычисляется по исходной $f$, вторая --- по касательной модели.
\end{frame}

% Slide 29
\begin{frame}{Условие Армихо: определение}
\small
\begin{block}{Определение}
При $0<c<1$ шаг $\alpha>0$ удовлетворяет \textbf{условию Армихо}, если
\[
\boxed{f(x+\alpha d)\le f(x)+c\alpha\nabla f(x)^\top d.}
\]
\end{block}
Эквивалентно:
\[
\boxed{\Delta_{\text{факт}}(\alpha)\ge c\,\Delta_{\text{лин}}(\alpha).}
\]
То есть реальное уменьшение не меньше доли $c$ от уменьшения по линейному приближению.
\end{frame}

% Slide 30
\begin{frame}{Геометрический смысл Армихо}
\small
\figtall{armijo_geometry.pdf}
\vspace{-.5em}
Принимаем те $\alpha$, для которых график $\varphi(\alpha)$ лежит ниже границы $\varphi(0)+c\alpha\varphi'(0)$.
\end{frame}

% Slide 31
\begin{frame}{Как на условие влияет параметр $c$}
\small
\figtall{armijo_c.pdf}
\vspace{-.4em}
Чем ближе $c$ к единице, тем ближе требование к полному линейному уменьшению; чем меньше $c$, тем оно мягче.
\end{frame}

% Slide 32
\begin{frame}{Почему недостаточно просто $f(x+\alpha d)<f(x)$}
\small
При любом направлении убывания очень маленький $\alpha$ уменьшит функцию, но может почти не изменить значение $f$.\\[.5em]
Условие Армихо требует \textbf{количественно заметного} убывания:
\[
f(x)-f(x+\alpha d)
\ge -c\alpha\nabla f(x)^\top d.
\]
При градиентном направлении правая часть равна
\[
\boxed{c\alpha\|\nabla f(x)\|_2^2.}
\]
Это не просто качественное условие уменьшения.
\end{frame}

% Slide 33
\begin{frame}{Арифметика Армихо на функции $f(x)=x^2$}
\small
Пусть $x=2$, $g=f'(2)=4$, $d=-4$. Тогда
\[
f(x+\alpha d)=(2-4\alpha)^2,\qquad
f(2)+c\alpha f'(2)d=4-16c\alpha.
\]
Условие Армихо:
\[
4-16\alpha+16\alpha^2\le4-16c\alpha.
\]
Для $\alpha>0$:
\[
\boxed{\alpha\le1-c.}
\]
При $c=\frac14$ подходят все $0<\alpha\le\frac34$.
\end{frame}

% Slide 34
\begin{frame}{Почему подходящий маленький шаг существует}
\small
\begin{block}{Теорема 1}
Если $f$ дифференцируема в точке $x$, а $d$ --- направление убывания, то при любом $c\in(0,1)$ \textbf{все достаточно малые} $\alpha>0$ удовлетворяют Армихо.
\end{block}
\textbf{Доказательство.} При $\alpha\to0+$
\[
f(x+\alpha d)=f(x)+\alpha\nabla f(x)^\top d+o(\alpha).
\]
Остаётся сравнить линейный член и остаток.
\end{frame}

% Slide 35
\begin{frame}{Доказательство существования шага Армихо}
\small
Обозначим $s=\nabla f(x)^\top d<0$.
Требуемое неравенство после подстановки разложения:
\[
\alpha s+o(\alpha)\le c\alpha s
\quad\Longleftrightarrow\quad
(1-c)s+\frac{o(\alpha)}{\alpha}\le0.
\]
Так как $(1-c)s<0$ и $o(\alpha)/\alpha\to0$, последнее неравенство выполняется при всех достаточно малых $\alpha>0$. $\square$
\end{frame}

% Slide 36
\begin{frame}{Напомним лемму об убывании}
\small
Если $f$ имеет $L$-липшицев градиент, то для любых $x,h$
\[
\boxed{f(x+h)\le f(x)+\nabla f(x)^\top h+\frac L2\|h\|_2^2.}
\]
Подставляя $h=\alpha d$, получаем
\[
f(x+\alpha d)\le f(x)+\alpha\nabla f(x)^\top d
+\frac{L\alpha^2}{2}\|d\|_2^2.
\]
Теперь можно вывести явную \textbf{достаточную границу} на шаг Армихо.
\end{frame}

% Slide 37
\begin{frame}{Достаточное условие Армихо при $L$-гладкости}
\small
Чтобы обеспечить Армихо, достаточно выполнения
\[
\alpha\nabla f(x)^\top d+\frac{L\alpha^2}{2}\|d\|_2^2
\le c\alpha\nabla f(x)^\top d.
\]
При $\alpha>0$, $L>0$ преобразуем:
\[
\frac{L\alpha}{2}\|d\|_2^2
\le (1-c)(-\nabla f(x)^\top d).
\]
Отсюда
\[
\boxed{0<\alpha\le \frac{2(1-c)(-\nabla f(x)^\top d)}{L\|d\|_2^2}.}
\]
\end{frame}

% Slide 38
\begin{frame}{Специальный случай: градиентное направление}
\small
При $d=-\nabla f(x)=-g$, $g\ne0$,
\[
-\nabla f(x)^\top d=\|g\|_2^2,\qquad \|d\|_2^2=\|g\|_2^2.
\]
Поэтому предыдущая оценка упрощается:
\[
\boxed{0<\alpha\le\bar\alpha:=\frac{2(1-c)}{L}.}
\]
\begin{alertblock}{Важно}
Это \textbf{достаточное}, но не обязательно необходимое условие Армихо. Более крупные шаги иногда тоже подходят.
\end{alertblock}
\end{frame}

% Slide 39
\begin{frame}{Как угол направления влияет на допустимый шаг}
\small
Пусть $\theta$ --- угол между $d$ и $-g$, $g=\nabla f(x)$.
Тогда $-g^\top d=\|g\|\|d\|\cos\theta$ и достаточная граница принимает вид
\[
\boxed{\alpha\le\frac{2(1-c)\|g\|\cos\theta}{L\|d\|}.}
\]
\figshort{angle_armijo.pdf}
\end{frame}

% Slide 40
\begin{frame}{Проверка оценок на квадратичном примере}
\small
Возьмём снова $f(x)=x^2$, $x=2$, $c=\frac14$.
Здесь глобальная константа гладкости $L=2$.
\[
\bar\alpha=\frac{2(1-c)}{L}=
\frac{2\cdot3/4}{2}=\frac34.
\]
Ранее прямым вычислением мы получили \textbf{точно тот же} интервал принятия:
\[
0<\alpha\le\frac34.
\]
Для других функций оценка через $L$ может оказаться консервативной.
\end{frame}

% Slide 41
\begin{frame}{Поиск шага с последовательным сокращением}
\small
Берём начальный пробный шаг $\alpha_0>0$ и коэффициент $\beta\in(0,1)$.\\[.5em]
Последовательно проверяем
\[
\boxed{\alpha_j=\alpha_0\beta^j,\qquad j=0,1,2,\dots}
\]
до первого выполнения условия Армихо.
\begin{block}{Терминология}
Этот алгоритм также называют \textbf{обратным поиском шага}. \textit{Обратным} здесь означает только последовательное \textbf{сокращение длины шага}, а не смену направления.
\end{block}
\end{frame}

% Slide 42
\begin{frame}{Алгоритм в псевдокоде}
\small
\begin{block}{Поиск шага с последовательным сокращением}
\begin{enumerate}
\item Вычислить $f(x)$, $g=\nabla f(x)$, выбрать $d$ и $\alpha\leftarrow\alpha_0$.
\item Пока 
\[
f(x+\alpha d)>f(x)+c\alpha g^\top d,
\]
заменить $\alpha\leftarrow\beta\alpha$.
\item Вернуть найденный $\alpha$.
\end{enumerate}
\end{block}
При $d=-g$ правая часть проверки: $f(x)-c\alpha\|g\|_2^2$.
\end{frame}

% Slide 43
\begin{frame}{Пример двух проверок: $f(x)=x^2$}
\small
\figtall{backtrack_example.pdf}
\vspace{-.35em}
При $x=2$, $\alpha_0=1$, $\beta=1/2$, $c=1/4$: шаг $1$ отклоняется, следующий шаг $1/2$ принимается.
\end{frame}

% Slide 44
\begin{frame}{Поиск шага всегда завершается}
\small
\begin{block}{Теорема 2}
Если $f$ дифференцируема в $x$, $d$ --- направление убывания, $\alpha_0>0$, $0<\beta,c<1$, то алгоритм завершится \textbf{за конечное число проверок}.
\end{block}
\textbf{Доказательство.} По теореме 1 существует $\delta>0$: все шаги $0<\alpha\le\delta$ удовлетворяют Армихо. Поскольку
\[
\alpha_0\beta^j\xrightarrow[j\to\infty]{}0,
\]
то для некоторого конечного $j$ имеем $\alpha_0\beta^j\le\delta$. Значит, поиск остановится. $\square$
\end{frame}

% Slide 45
\begin{frame}{Сколько проверок может понадобиться?}
\small
Если функция $L$-гладкая и $d=-g$, то достаточно
\[
\alpha_0\beta^j\le\bar\alpha=\frac{2(1-c)}{L}.
\]
При $\alpha_0>\bar\alpha$ это выполнится, когда
\[
j\ge\frac{\ln(\bar\alpha/\alpha_0)}{\ln\beta},
\qquad\ln\beta<0.
\]
Следовательно, число \textbf{проверок функции} не превышает
\[
\boxed{1+\max\!\left\{0,\left\lceil\frac{\ln(\bar\alpha/\alpha_0)}{\ln\beta}\right\rceil\right\}.}
\]
\end{frame}

% Slide 46
\begin{frame}{Нижняя граница на принятый шаг}
\small
\begin{block}{Утверждение 5}
Для $L$-гладкой функции при $d=-g$, $g\ne0$, поиск с фиксированными $\alpha_0,\beta,c$ возвращает $\alpha_k$, удовлетворяющий
\[
\boxed{\alpha_k\ge\alpha_{\min}:=\min\!\left\{\alpha_0,\,\beta\frac{2(1-c)}{L}\right\}>0.}
\]
\end{block}
Эта оценка важна: подбор шага \textbf{не допускает его бесконтрольного уменьшения} при заданных предпосылках.
\end{frame}

% Slide 47
\begin{frame}{Доказательство нижней границы}
\small
Обозначим $\bar\alpha=2(1-c)/L$.
\begin{itemize}
\item Если приняли $\alpha_0$, то $\alpha_k=\alpha_0\ge\alpha_{\min}$.
\item Если $\alpha_k=\beta\alpha_{\mathrm{prev}}$ получен после хотя бы одного отказа, то предыдущий шаг $\alpha_{\mathrm{prev}}$ не удовлетворял Армихо.
\end{itemize}
Но \textbf{все} шаги $\alpha\le\bar\alpha$ его удовлетворяют. Значит,
\[
\alpha_{\mathrm{prev}}>\bar\alpha
\quad\Longrightarrow\quad
\alpha_k=\beta\alpha_{\mathrm{prev}}>\beta\bar\alpha.
\]
В обоих случаях $\alpha_k\ge\alpha_{\min}$. $\square$
\end{frame}

% Slide 48
\begin{frame}{Как влияет коэффициент сокращения $\beta$}
\small
\figtall{beta_decay.pdf}
\vspace{-.4em}
При большом $\beta$ проверки чаще дают шаги близких величин. Малое $\beta$ быстрее уменьшает шаг, но может привести к излишне короткому принятому шагу.
\end{frame}

% Slide 49
\begin{frame}{Сравнение стратегий на одной задаче}
\small
\figtall{strategy_comparison.pdf}
\vspace{-.45em}
Фиксированный, точный и адаптивно сокращаемый шаг могут существенно отличаться как по числу итераций, так и по числу вычислений $f$.
\end{frame}

% Slide 50
\begin{frame}{Цена одного шага алгоритма}
\small
Пусть вычисление $f(x)$ стоит $C_f$, а $\nabla f(x)$ --- $C_g$.\\[.4em]
\begin{itemize}
\item \textbf{Фиксированный шаг:} обычно одна оценка градиента и одно новое значение функции.
\item \textbf{Точный поиск:} дополнительная одномерная задача оптимизации.
\item \textbf{Поиск с сокращением:} одна оценка $\nabla f(x)$, затем несколько проверок $f(x+\alpha d)$.
\end{itemize}
\begin{block}{Вывод}
Меньшее число итераций не обязательно означает меньшую вычислительную стоимость.
\end{block}
\end{frame}

% Slide 51
\begin{frame}{Гарантия уменьшения на каждом шаге}
\small
Для градиентного направления принятый по Армихо шаг удовлетворяет
\[
f(x_{k+1})\le f(x_k)-c\alpha_k\|\nabla f(x_k)\|_2^2.
\]
По доказанной нижней границе $\alpha_k\ge\alpha_{\min}>0$, поэтому
\[
\boxed{f(x_{k+1})\le f(x_k)-c\alpha_{\min}\|\nabla f(x_k)\|_2^2.}
\]
Теперь можно получить количественную оценку на стационарность без выпуклости.
\end{frame}

% Slide 52
\begin{frame}{Суммируем убывания по итерациям}
\small
Просуммируем от $k=0$ до $N-1$:
\[
\begin{aligned}
c\alpha_{\min}\sum_{k=0}^{N-1}\|\nabla f(x_k)\|_2^2
&\le\sum_{k=0}^{N-1}\big(f(x_k)-f(x_{k+1})\big)\\
&=f(x_0)-f(x_N).
\end{aligned}
\]
Если $f(x)\ge f_{\inf}$ для всех $x$, то
\[
\boxed{\sum_{k=0}^{N-1}\|\nabla f(x_k)\|_2^2
\le\frac{f(x_0)-f_{\inf}}{c\alpha_{\min}}.}
\]
\end{frame}

% Slide 53
\begin{frame}{Следствие: градиент стремится к нулю}
\small
\begin{block}{Следствие из оценки суммы}
Если $f$ $L$-гладкая, ограничена снизу, а на каждом шаге применяется поиск с фиксированными $\alpha_0,\beta,c$ и направлением $-\nabla f(x_k)$, то
\[
\boxed{\|\nabla f(x_k)\|_2\longrightarrow0.}
\]
\end{block}
\textbf{Почему?} Неотрицательный ряд
$\sum_{k\ge0}\|\nabla f(x_k)\|_2^2$
имеет ограниченные частичные суммы, значит сходится, а его члены стремятся к нулю.
\end{frame}

% Slide 54
\begin{frame}{Количественная оценка: $O(1/\sqrt N)$}
\small
Из оценки суммы дополнительно получаем
\[
\min_{0\le k<N}\|\nabla f(x_k)\|_2^2
\le\frac{f(x_0)-f_{\inf}}{c\alpha_{\min}N}.
\]
Следовательно,
\[
\boxed{\min_{0\le k<N}\|\nabla f(x_k)\|_2
\le\sqrt{\frac{f(x_0)-f_{\inf}}{c\alpha_{\min}N}}.}
\]
Это гарантия достижения \textbf{приближённой стационарной точки}, а не глобального минимума.
\end{frame}

% Slide 55
\begin{frame}{Критерии остановки: что проверять?}
\small
\begin{block}{Три популярных правила}
\begin{enumerate}
\item $\|\nabla f(x_k)\|_2\le\varepsilon_g$;
\item $\|x_{k+1}-x_k\|_2\le\varepsilon_x$;
\item $|f(x_{k+1})-f(x_k)|\le\varepsilon_f$.
\end{enumerate}
\end{block}
Первое правило измеряет приближённую стационарность. Два других могут сработать \textbf{из-за слишком короткого шага}, даже вдали от минимума.
\end{frame}

% Slide 56
\begin{frame}{Почему маленькое изменение — не признак оптимума}
\small
\figtall{stopping_pitfalls.pdf}
\vspace{-.4em}
Слева малое $\|x_{k+1}-x_k\|$ может возникнуть из-за $\alpha_k\approx0$. Справа $\nabla f=0$ может означать седловую точку.
\end{frame}

% Slide 57
\begin{frame}{Итоги и следующая лекция}
\small
\begin{block}{Что мы доказали сегодня}
Точный поиск вдоль антиградиента даёт $g_{k+1}^\top g_k=0$.
Поиск с последовательным сокращением и условием Армихо при
$L$-гладкости обеспечивает положительную нижнюю границу шага.
Если $f$ ограничена снизу, то $\|\nabla f(x_k)\|_2\to0$.
\end{block}

\begin{alertblock}{Лекция 9. Метод Ньютона}
Сегодня мы подбирали \textbf{длину шага вдоль заданного направления}.
А что если выбирать само перемещение, минимизируя
\textbf{квадратичное приближение} функции?
\end{alertblock}
\[
m_x(h)=f(x)+\nabla f(x)^\top h+
\frac12 h^\top\nabla^2f(x)h.
\]
При $\nabla^2f(x)\succ0$ минимум этой модели достигается при
\[
\boxed{h_{\mathrm N}=-\bigl(\nabla^2f(x)\bigr)^{-1}\nabla f(x).}
\]
\end{frame}

\end{document}
